BGV supports exact arithmetic in its finite plaintext ring, subject to decryption correctness. BFV uses a large message scale $\Delta$, while BGV places the error in a multiple of its plaintext modulus. The familiar high/low digit picture is an intuition; a literal split into binary bit positions applies only with suitable power-of-two bases, and centered representatives can be negative.

BGV shares exact plaintext-ring arithmetic, tensor multiplication, and key-switching techniques with BFV. BGV commonly uses modulus switching to control noise magnitude along a modulus chain. Fixed-modulus BFV is also limited by its noise budget; it does not support an unlimited number of multiplications without refreshing. Modulus switching and bootstrapping require the scheme-specific correctness conditions described below.



\begin{tcolorbox}[
    title = \textbf{Required Background},    % box title
    colback = white,    % light background; tweak to taste
    colframe = black,  % frame colour
    boxrule = 0.8pt,     % line thickness
    left = 1mm, right = 1mm, top = 1mm, bottom = 1mm % inner padding
]

\begin{itemize}
\item \autoref{sec:modulo}: \nameref{sec:modulo}
\item \autoref{sec:group}: \nameref{sec:group}
\item \autoref{sec:field}: \nameref{sec:field}
\item \autoref{sec:order}: \nameref{sec:order}
\item \autoref{sec:polynomial-ring}: \nameref{sec:polynomial-ring}
\item \autoref{sec:decomp}: \nameref{sec:decomp}
\item \autoref{sec:roots}: \nameref{sec:roots}
\item \autoref{sec:cyclotomic}: \nameref{sec:cyclotomic}
\item \autoref{sec:cyclotomic-polynomial-integer-ring}: \nameref{sec:cyclotomic-polynomial-integer-ring}
\item \autoref{sec:matrix}: \nameref{sec:matrix}
\item \autoref{sec:euler}: \nameref{sec:euler}
\item \autoref{sec:modulus-rescaling}: \nameref{sec:modulus-rescaling}
\item \autoref{sec:chinese-remainder}: \nameref{sec:chinese-remainder}
\item \autoref{sec:polynomial-interpolation}: \nameref{sec:polynomial-interpolation}
\item \autoref{sec:ntt}: \nameref{sec:ntt}
\item \autoref{sec:lattice}: \nameref{sec:lattice}
\item \autoref{sec:rlwe}: \nameref{sec:rlwe}
\item \autoref{sec:glwe}: \nameref{sec:glwe}
\item \autoref{sec:glwe-add-cipher}: \nameref{sec:glwe-add-cipher}
\item \autoref{sec:glwe-add-plain}: \nameref{sec:glwe-add-plain}
\item \autoref{sec:glwe-mult-plain}: \nameref{sec:glwe-mult-plain}
\item \autoref{subsec:modulus-switch-rlwe}: \nameref{subsec:modulus-switch-rlwe}
\item \autoref{sec:glwe-key-switching}: \nameref{sec:glwe-key-switching}
\item \autoref{sec:bfv}: \nameref{sec:bfv}
\item \autoref{sec:ckks}: \nameref{sec:ckks}
\end{itemize}
\end{tcolorbox}

\clearpage

\subsection{Encoding and Decoding}
\label{subsec:bgv-encoding-decoding}

BGV uses almost the same plaintext encoding scheme as BFV as described in Summary~\ref{subsubsec:bfv-encoding-summary} in \autoref{subsubsec:bfv-encoding-summary}, with the only difference that the scaling factor $\Delta = \left\lfloor\dfrac{q}{t}\right\rfloor$ is not applied to the plaintext polynomial $M(X)$ like BFV does. Instead, BGV applies its own scaling factor $\Delta = t$ to the noise polynomial $E(X)$ whenever it encrypts a new ciphertext (will be explained in \autoref{subsec:bgv-enc-dec}).

The following is BGV's encoding and decoding scheme.

\begin{tcolorbox}[title={\textbf{\tboxlabel{\ref*{subsec:bgv-encoding-decoding}} BGV's Encoding and Decoding}}]



\textbf{\underline{Input}:} An $n$-slot integer modulo $t$ vector $\vec{v} = (v_0, v_1, \cdots, v_{n-1}) \in \mathbb{Z}_t^n$

$ $

\textbf{\underline{Encoding}:}

 Convert $\vec{v} \in \mathbb{Z}_t^n$ into $\vec{m} \in \mathbb{Z}_t^n$ by applying the transformation $\vec{m} = \dfrac{\hathat W \cdot I_n^R \cdot \vec{v}}{n}$

, where $\hathat W$ is a basis of the $n$-dimensional vector space crafted as follows:

{\footnotesize{$\hathat W = \begin{bmatrix}
1 & 1 & \cdots & 1 & 1 & 1 & \cdots & 1\\
(\omega^{J(\frac{n}{2} - 1)}) & (\omega^{J(\frac{n}{2} - 2)}) & \cdots & (\omega^{J(0)}) & (\omega^{J_*(\frac{n}{2} - 1)}) & (\omega^{J_*(\frac{n}{2} - 2)}) & \cdots & (\omega^{J_*(0)})\\
(\omega^{J(\frac{n}{2} - 1)})^2 & (\omega^{J(\frac{n}{2} - 2)})^2 & \cdots & (\omega^{J(0)})^2 & (\omega^{J_*(\frac{n}{2} - 1)})^2 & (\omega^{J_*(\frac{n}{2} - 2)})^2 & \cdots & (\omega^{J_*(0)})^2 \\
\vdots & \vdots & \ddots & \vdots & \vdots & \ddots & \vdots & \vdots \\
(\omega^{J(\frac{n}{2} - 1)})^{n-1} & (\omega^{J(\frac{n}{2} - 2)})^{n-1} & \cdots & (\omega^{J(0)})^{n-1} & (\omega^{J_*(\frac{n}{2} - 1)})^{n-1} & (\omega^{J_*(\frac{n}{2} - 2)})^{n-1} & \vdots  & (\omega^{J_*(0)})^{n-1}
\end{bmatrix}$}}

$ $

, where for this fully split batching construction $t$ is prime and $t\equiv1\pmod{2n}$. Choose a generator $g$ of $\mathbb{Z}_t^\times$ and $\omega=g^{(t-1)/(2n)}\bmod t$ (\autoref{subsubsec:poly-vector-transformation-modulus}). The output is $M=\sum_{i=0}^{n-1}m_iX^i\in\mathbb{Z}_t[X]/(X^n+1)$. Its coefficients are lifted to integer representatives for encryption modulo $q$; this lift is not an inclusion of rings. More general BGV plaintext moduli do not necessarily provide $n$ scalar slots.

Choose a fixed representative of each coefficient modulo $t$ before embedding it into the ciphertext coefficient domain. Correctness is determined by the full phase bound, not just by the sizes of the fresh message coefficients.

$ $

\textbf{\underline{Decoding}:} For the plaintext polynomial $M = \sum\limits_{i=0}^{n-1} m_iX^i$,
compute $\vec{v} = \hathat W^* \cdot \vec{m}$, where

$\hathat{W}^* = \begin{bmatrix}
1 & (\omega^{J(0)}) & (\omega^{J(0)})^2 & \cdots & (\omega^{J(0)})^{n-1}\\
1 & (\omega^{J(1)}) & (\omega^{J(1)})^2 & \cdots & (\omega^{J(1)})^{n-1}\\
1 & (\omega^{J(2)}) & (\omega^{J(2)})^2 & \cdots & (\omega^{J(2)})^{n-1}\\
\vdots & \vdots & \vdots & \ddots & \vdots \\
1 & (\omega^{J(\frac{n}{2}-1)}) & (\omega^{J(\frac{n}{2}-1)})^2 & \cdots & (\omega^{J(\frac{n}{2}-1)})^{n-1}\\
1 & (\omega^{J_*(0)}) & (\omega^{J_*(0)})^2 & \cdots & (\omega^{J_*(0)})^{n-1}\\
1 & (\omega^{J_*(1)}) & (\omega^{J_*(1)})^2 & \cdots & (\omega^{J_*(1)})^{n-1}\\
1 & (\omega^{J_*(2)}) & (\omega^{J_*(2)})^2 & \cdots & (\omega^{J_*(2)})^{n-1}\\
\vdots & \vdots & \vdots & \ddots & \vdots \\
1 & (\omega^{J_*(\frac{n}{2}-1)}) & (\omega^{J_*(\frac{n}{2}-1)})^2 & \cdots & (\omega^{J_*(\frac{n}{2}-1)})^{n-1}\\
\end{bmatrix}$

\end{tcolorbox}


\subsection{Encryption and Decryption}
\label{subsec:bgv-enc-dec}

Compare the phase $\Delta M+E$ of BFV, with $\Delta=\lfloor q/t\rfloor$, to the phase $M+tE$ of BGV (\autoref{subsec:bfv-enc-dec}). In BGV, reduction modulo the plaintext modulus removes the divisible error after the centered phase has been recovered correctly. The following symmetric-key presentation uses a prime plaintext modulus and a convenient congruent modulus chain; BGV also has more general plaintext and chain configurations.

Also, in BGV, the ciphertext modulus $q$ is leveled like CKKS's one: $q \in \{q_0, q_1, \cdots, q_L\}$, where each $q_l = \prod\limits_{i=0}^l w_i$ (where each $w_i$ is a CRT modulus).

BGV's encryption decryption process is described as follows:


\begin{tcolorbox}[title={\textbf{\tboxlabel{\ref*{subsec:bgv-enc-dec}} BGV Encryption and Decryption}}]

\textbf{\underline{Initial Setup}:}

\begin{itemize}
\item The plaintext modulus $t = p$ (a prime)

\item The ciphertext modulus $q$ is leveled like in CKKS: $q \in \{q_0, q_1, \cdots, q_L\}$, where each $q_l = \prod\limits_{i=0}^l w_i$ (each $w_i$ is a CRT modulus), and each $q_l \equiv 1 \bmod t$ (will be explained in \autoref{subsec:bgv-modulus-switch})

\item The noise scaling factor $\Delta=t$
\item The secret key $S \xleftarrow{\$} \mathcal{R}_{\langle n, \textit{tern} \rangle}$. The coefficients of the polynomial $S$ ternary (i.e., $\{-1, 0, 1\}$).

\end{itemize}

\par\noindent\rule{\textwidth}{0.4pt}

\textbf{\underline{Encryption Input}:} $M \in \mathcal{R}_{\langle n, q \rangle}$, $A \xleftarrow{\$} \mathcal{R}_{\langle n, q \rangle}$, $E \xleftarrow{\chi_\sigma} \mathcal{R}_{\langle n, q \rangle}$


\begin{enumerate}
%\item Scale up $M \rightarrow \Delta M \text { } \in \mathcal{R}_{\langle n, q\rangle}$

\item Compute $B = -A \cdot S + M + \Delta E \text{ } \in \mathcal{R}_{\langle n,q \rangle}$

\item $\textsf{RLWE}_{S,\sigma}(M + \Delta E) = (A, B) \text{ } \in \mathcal{R}_{\langle n,q \rangle}^2$

\end{enumerate}

\par\noindent\rule{\textwidth}{0.4pt}

\textbf{\underline{Decryption Input}:} $\textsf{ct} = (A, B) \text{ } \in \mathcal{R}_{\langle n,q \rangle}^2$

\begin{enumerate}
\item $\textsf{RLWE}^{-1}_{S,\sigma}(\textsf{ct}) = B + A \cdot S  = M + \Delta E \pmod q$

\item Compute $M + \Delta E \bmod t$ to get $M$. \textcolor{red}{ $\rhd$ modulo reduction of $M + \Delta E$ by $t$}

\end{enumerate}

$ $

The final output is $M(X) \in \mathbb{Z}_t[X] / (X^n + 1)$.

$ $

\textbf{\underline{Conditions for Correct Decryption}:}

With centered representatives, require $\|M+tE\|_\infty<q_l/2$ at the level where decryption is performed. First compute $[B+AS]_{q_l}$ in the centered interval and then reduce modulo $t$. The one-sided inequality $m_i+te_i<q_l$ alone is insufficient, and reducing modulo $q_l$ is not generally compatible with subsequent reduction modulo $t$.

%\end{enumerate}

\end{tcolorbox}

The term eliminated by reduction modulo $t$ is $tE$, whose coefficients are multiples of $t$; the coefficients of $E$ need not be multiples of $t$. The same rule must hold for evaluation-key errors in key switching and relinearization. Throughout this chapter, encryption notation records the phase relation and does not assert that derived errors retain the original distribution.



\subsection{Ciphertext-to-Ciphertext Addition}
\label{subsec:bgv-add-cipher}

BGV's ciphertext-to-ciphertext addition scheme is exactly the same as BFV's scheme (Summary~\ref*{subsec:bfv-add-cipher} in \autoref{subsec:bfv-add-cipher}).


\begin{tcolorbox}[title={\textbf{\tboxlabel{\ref*{subsec:bgv-add-cipher}} BGV Ciphertext-to-Ciphertext Addition}}]
$\textsf{RLWE}_{S, \sigma}(M^{\langle 1 \rangle} + \Delta E^{\langle 1 \rangle} ) + \textsf{RLWE}_{S, \sigma}(M^{\langle 2 \rangle} + \Delta E^{\langle 2 \rangle}) $

$ = ( A^{\langle 1 \rangle}, \text{ } B^{\langle 1 \rangle}) + (A^{\langle 2 \rangle}, \text{ } B^{\langle 2 \rangle}) $

$ = ( A^{\langle 1 \rangle} + A^{\langle 2 \rangle}, \text{ } B^{\langle 1 \rangle} + B^{\langle 2 \rangle} ) $

$= \textsf{RLWE}_{S, \sigma}\bm((M^{\langle 1 \rangle} + M^{\langle 2 \rangle})  + \Delta E^{\langle 1 \rangle} + \Delta E^{\langle 2 \rangle} \bm)$
\end{tcolorbox}

\subsection{Ciphertext-to-Plaintext Addition}
\label{subsec:bgv-add-plain}

BGV's ciphertext-to-plaintext addition scheme is almost the same as BFV's scheme (Summary~\ref*{subsec:bfv-add-plain} in \autoref{subsec:bfv-add-plain}). However, one difference is that it's not the case that the plaintext polynomial $\Lambda(X)$ to be added is scaled up by $\Delta$, but it remains as $\Lambda(X)$.

\begin{tcolorbox}[title={\textbf{\tboxlabel{\ref*{subsec:bgv-add-plain}} BGV Ciphertext-to-Plaintext Addition}}]
$\textsf{RLWE}_{S, \sigma}(M + \Delta E) + \Lambda $

$=  (A, \text{ } B) + \Lambda$

$=  (A, \text{ } B + \Lambda)$

$= \textsf{RLWE}_{S, \sigma}\bm((M + \Lambda) + \Delta E \bm)$
\end{tcolorbox}


\subsection{Ciphertext-to-Plaintext Multiplication}
\label{subsec:bgv-mult-plain}




BGV's ciphertext-to-plaintext multiplication scheme is exactly the same as BFV's scheme (Summary~\ref*{subsec:bfv-mult-plain} in \autoref{subsec:bfv-mult-plain}).


\begin{tcolorbox}[title={\textbf{\tboxlabel{\ref*{subsec:bgv-mult-plain}} BGV Ciphertext-to-Plaintext Multiplication}}]
$\textsf{RLWE}_{S, \sigma}(M + \Delta E) \cdot \Lambda$

$= (A, \text{ } B) \cdot \Lambda$

$= (A \cdot \Lambda, \text{ }  B \cdot \Lambda )$

$= \textsf{RLWE}_{S, \sigma}((M \cdot \Lambda) + \Delta E\cdot \Lambda )$
\end{tcolorbox}

Notice that BGV's ciphertext-to-plaintext multiplication does not consume any multiplicative level.

\subsection{\textsf{ModDrop}}
\label{subsec:bgv-moddrop}

BGV ModDrop reduces $(A,B)$ coefficient-wise modulo a divisor $q_{l-1}$ of $q_l$:
\[
(A',B')=([A]_{q_{l-1}},[B]_{q_{l-1}}).
\]
Since $A'S+B'\equiv AS+B\pmod{q_{l-1}}$, it preserves the centered phase $M+tE$ if $\|M+tE\|_\infty<q_{l-1}/2$. Under that condition the plaintext and noise remain unchanged while the modulus headroom decreases. Without the bound, wraparound can alter the plaintext modulo $t$.

\subsection{Modulus Switch}
\label{subsec:bgv-modulus-switch}

\noindent \textbf{- References:} \href{https://eprint.iacr.org/2020/1481}{Design and implementation of HElib}~\cite{bgv-modswitch1}; \href{https://eprint.iacr.org/2011/277.pdf}{(Leveled) Fully Homomorphic Encryption without Bootstrapping}~\cite{bgv-modswitch2}; \href{https://eprint.iacr.org/2012/099.pdf}{Homomorphic Evaluation of the AES Circuit}~\cite{bgv-modswitch3}.

We describe a convenient special case with $q_l\equiv\widehat q\equiv1\pmod t$ and $\widehat q<q_l$. These congruences are sufficient for this presentation, not necessary conditions for every BGV implementation. Let $\rho=\widehat q/q_l$ and take centered integer representatives of $A,B$. Compute the integer polynomials
\[
A'=\lceil\rho A\rfloor,\quad B'=\lceil\rho B\rfloor,\quad
\epsilon'_A=\widehat qA-q_lA',\quad\epsilon'_B=\widehat qB-q_lB'.
\]
The errors $\epsilon'_A/q_l,\epsilon'_B/q_l$ are real rounding errors of magnitude at most $1/2$, not elements of an integer residue ring. Correct the residues modulo $t$ by
\[
H_A=[q_l^{-1}\epsilon'_A]_t,\quad H_B=[q_l^{-1}\epsilon'_B]_t,
\quad (\widehat A,\widehat B)=([A'+H_A]_{\widehat q},[B'+H_B]_{\widehat q}).
\]
Here $H_A,H_B$ use centered representatives. Before the final reduction, $\widetilde A=A'+H_A$ and $\widetilde B=B'+H_B$ satisfy $\widetilde A\equiv A\pmod t$ and $\widetilde B\equiv B\pmod t$.

For a proof, write $AS+B=M+tE+Kq_l$ and set $\delta_A=\widetilde A-\rho A$, $\delta_B=\widetilde B-\rho B$. Then
\[
Y=\widetilde AS+\widetilde B-K\widehat q
=\rho(M+tE)+\delta_AS+\delta_B,
\qquad Y\equiv M\pmod t.
\]
The latter congruence follows from $q_l\equiv\widehat q\pmod t$. If $\|Y\|_\infty<\widehat q/2$, then $Y=[\widehat AS+\widehat B]_{\widehat q}$ and decryption returns $M$. A sufficient deterministic condition is
\[
\rho\|M+tE\|_\infty+\frac{t+1}{2}(1+\|S\|_1)<\frac{\widehat q}{2}.
\]
This follows from $\|\delta_A\|_\infty,\|\delta_B\|_\infty\leq(t+1)/2$. The bounded corrections are not automatically harmless when the old phase is close to its boundary.

\begin{tcolorbox}[title={\textbf{\tboxlabel{\ref*{subsec:bgv-modulus-switch}} BGV Modulus Switch}}]
Scale ciphertext representatives approximately by $\rho$, then apply the congruence correction above. Under the phase bound, the result has phase $M+tE'$ for a new integer error $E'$. The plaintext modulus and noise divisibility factor remain $t$; the magnitude of the old phase is reduced, with a new rounding/correction term. This is not a noise reset.
\end{tcolorbox}

\subsubsection{Difference between Modulus Switch and \textsf{ModDrop}}
\label{subsubsec:bgv-moddrop-vs-modswitch}

CKKS rescaling divides its message scale by $q_l/q_{l-1}$; its ModDrop leaves the scale unchanged under a phase bound (\autoref{subsubsec:ckks-mult-cipher-rescale}). In BGV, the noise divisibility factor $\Delta=t$ is unchanged by either operation. BGV modulus switching approximately reduces the magnitude of the phase and preserves the plaintext residue through a congruence correction. ModDrop performs only residue reduction and reduces the available phase headroom. Neither operation is restricted to the end of a computation.

\para{Source Code:} The \href{https://github.com/wang2000star/fhe-textbook-cn/blob/chinese-edition/source%20code/bgv.py}{\underline{Python example}} illustrates modulus switching; any implementation must enforce the phase and parameter conditions above.


\subsection{Ciphertext-to-Ciphertext Multiplication}
\label{subsec:bgv-mult-cipher}

For input phases $M_1+tE_1$ and $M_2+tE_2$, tensoring produces
\[
D_0=B_1B_2,\quad D_1=A_1B_2+A_2B_1,\quad D_2=A_1A_2,
\]
\[
D_0+D_1S+D_2S^2\equiv M_1M_2+t(M_1E_2+M_2E_1+tE_1E_2)\pmod q.
\]
The noise remains divisible by $t$: it is incorrect to regard the scheme's factor $t$ as successively becoming $t^2,t^4,\ldots$. Reducing $M_1M_2$ modulo $t$ can add an additional carry to the integer error.

\begin{tcolorbox}[title={\textbf{\tboxlabel{\ref*{subsec:bgv-mult-cipher}} BGV Ciphertext-to-Ciphertext Multiplication}}]
\begin{enumerate}
\item Compute $(D_0,D_1,D_2)$ modulo $q$.
\item Choose an exact integer gadget $g_i$ with $D_2\equiv\sum_i d_ig_i\pmod q$. The BGV evaluation keys must satisfy $B_i+A_iS\equiv g_iS^2+tE_i\pmod q$. Form $(A,B)=(D_1,D_0)+\sum_i d_i(A_i,B_i)$; its added key-switching error is $t\sum_i d_iE_i$.
\item If needed for noise management, apply the BGV modulus switch in \autoref{subsec:bgv-modulus-switch}. This uses the actual ratio $\widehat q/q$ and a congruence correction. It is not division by $t$ with rounding to a multiple of $t$.
\end{enumerate}
\para{Swapping the Order of \textsf{Relinearization} and \textsf{Rescaling}:} One may modulus-switch the three components first, if the correctness bound includes the secret vector $(1,S,S^2)$ and the relinearization key is valid at the new modulus. The order changes the error bound and cost.
\end{tcolorbox}

\para{Details of the Optional Rescaling:} The new error depends on the input messages, both input errors, polynomial convolution, key-switching, and the modulus-switch correction. Merely choosing $w_l=q_l/q_{l-1}>t$ does not guarantee that the output error is smaller than both input errors. Select the modulus chain using a complete bound on the resulting phase. CKKS likewise uses an explicit scale and error budget; noise growth is not a universal monotonic statement about every input.


\subsection{Homomorphic Key Switching}
\label{subsec:bgv-key-switching}

Choose an exact integer gadget $g_i$ and digits $d_i$ satisfying $A\equiv\sum_i d_ig_i\pmod q$. A BGV evaluation key under $S'$ must satisfy $b_i+a_iS'\equiv g_iS+tE_i\pmod q$. Then
\begin{tcolorbox}[title={\textbf{\tboxlabel{\ref*{subsec:bgv-key-switching}} BGV Key Switching}}]
\[
(A',B')=(0,B)+\sum_i d_i(a_i,b_i),\qquad
B'+A'S'\equiv M+t\left(E+\sum_i d_iE_i\right)\pmod q.
\]
\end{tcolorbox}
This is the gadget technique of \autoref{subsec:bfv-key-switching}, with the BGV-specific factor $t$ in the evaluation-key errors. Check the resulting centered phase bound. Using generic evaluation keys with unscaled errors can change the plaintext residue.

\subsection{Homomorphic Rotation of Input Vector Slots}
\label{subsec:bgv-rotation}

For fully split batching and the root order used in \autoref{subsec:bfv-rotation}, the automorphism $X\mapsto X^{J(h)}$, $J(h)=5^h\bmod2n$, rotates the two groups of $n/2$ scalar slots left by $h$ positions within each group.
\begin{tcolorbox}[title={\textbf{\tboxlabel{\ref*{subsec:bgv-rotation}} BGV Homomorphic Rotation of Input Vector Slots}}]
\begin{enumerate}
\item Replace $(A(X),B(X))$ by $(A(X^{J(h)}),B(X^{J(h)}))$ modulo $X^n+1$. The secret becomes $S(X^{J(h)})$ and the phase is $M(X^{J(h)})+tE(X^{J(h)})$ modulo $q$.
\item Apply \autoref{subsec:bgv-key-switching} to switch from $S(X^{J(h)})$ back to $S(X)$, using evaluation keys whose errors are multiples of $t$.
\end{enumerate}
\end{tcolorbox}
The automorphism permutes the message exactly, while key switching adds error. Correctness requires the resulting phase to satisfy its centered no-wrap bound. Other batching layouts have their corresponding automorphism-induced permutations.


\subsection{Modulus Bootstrapping}
\label{subsec:bgv-bootstrapping}

\noindent \textbf{- Reference:} \href{https://eprint.iacr.org/2022/1363}{Bootstrapping for BGV and BFV Revisited}~\cite{cryptoeprint:2022/1363}, Sections 3.2--3.4 and 6.

Bootstrapping evaluates a suitable decryption circuit using an encrypted secret key. It produces a ciphertext with a renewed modulus and noise budget, after paying for the operations of the circuit itself. The following is the general bootstrapping outline of the cited paper; it is not an assertion that merely lifting a BGV ciphertext to a larger modulus refreshes it.

The coefficient-to-slot route below also assumes $\gcd(p,\mu)=1$ for the cyclotomic index $\mu$, as in Section 2.3 of the cited paper. For this book's power-of-two ring with $\mu=2n$, this requires odd $p$. The paper's binary digit-extraction case uses compatible cyclotomic parameters; its rounding rule alone does not provide scalar or extension-ring slots in the ramified power-of-two ring.

Let the plaintext modulus be $t=p^r$, where $p$ is prime and $r\geq1$. After a valid modulus switch, let the ciphertext modulus be $q$ with $q\equiv1\pmod{p^\varepsilon}$, where $\varepsilon>r$. Write its centered decryption phase as $M+p^rE$. Choose centered coefficient representatives and compute publicly
\[
A'=[p^{\varepsilon-r}A]_q,\qquad B'=[p^{\varepsilon-r}B]_q.
\]
There is an integer polynomial $K'$ satisfying
\[
u=B'+A'S=p^{\varepsilon-r}M+p^\varepsilon E+qK'.
\]
Consequently,
\[
w=[u]_{p^\varepsilon}=[p^{\varepsilon-r}M+K']_{p^\varepsilon},\qquad
M=\left[\left\lfloor\frac{w}{p^{\varepsilon-r}}\right\rceil\right]_{p^r}
\]
provided that $\|K'\|_\infty<p^{\varepsilon-r}/2$. This is the rounded decryption formula of the cited paper, not an ordinary floor. A sufficient bound must account for all three terms
\[
\|K'\|_\infty\leq\left\|\frac{B'+A'S}{q}\right\|_\infty+
\frac{p^{\varepsilon-r}\|M\|_\infty}{q}+
\frac{p^\varepsilon\|E\|_\infty}{q}.
\]
In particular, the bound $\|K'\|_\infty\leq n+1$ does not follow merely from a ternary secret. The input noise and the chosen parameters must also be bounded.

\begin{enumerate}
\item \textbf{Homomorphic inner product.} Use a bootstrapping key encrypting $S$ under the output secret, at a sufficiently large ciphertext modulus $Q$ and plaintext modulus $p^\varepsilon$. Evaluate $B'+A'S$ homomorphically. This encrypts $w=[u]_{p^\varepsilon}$, with a new encryption error divisible by $p^\varepsilon$. The old ciphertext coefficients $A',B'$ are public inputs to this computation. Correctness requires that this new ciphertext and every later intermediate ciphertext satisfy their respective decryption bounds.
\item \textbf{\textsf{CoeffToSlot}.} Homomorphically apply the appropriate linear transformation over the plaintext ring to place the coefficients of $w$ into scalar slots. For the fully split case $p\equiv1\pmod{2n}$, the ring has $n$ scalar slots and the transform is the inverse of the encoding matrix over $\mathbb Z_{p^\varepsilon}$. For general cyclotomic parameters, slots can be extension rings and multiple ciphertexts are needed, as explained in the cited paper. A complex CKKS matrix cannot simply be reused over this ring.
\item \textbf{Digit extraction.} Evaluate the slot-wise rounded map $w_i\mapsto[\lfloor w_i/p^{\varepsilon-r}\rceil]_{p^r}$. For odd $p$, it can be organized as repeated removal of balanced least significant base-$p$ digits and exact division by $p$. The digit-extraction polynomials must return the required balanced digit modulo the current plaintext modulus; merely taking an ordinary floor fails for negative $K'$. Binary digit extraction needs its own rounding convention (including the offset in the cited paper). The input bound above, the digit-polynomial correctness conditions, and the homomorphic error budget are all required.
\item \textbf{\textsf{SlotToCoeff}.} Apply the inverse linear transformation to recover the polynomial $M$ in the plaintext ring modulo $p^r$. The output is a BGV ciphertext at the remaining ciphertext modulus $Q_{\mathrm{out}}$, with phase $M+p^rE_{\mathrm{out}}$. Its error is determined by the bootstrapping circuit; it is not the original error reset to zero.
\end{enumerate}

\para{\underline{Handling the Divisions by $p$}:} At an intermediate plaintext modulus $p^j$, suppose digit subtraction gives a phase $pT+p^jF$ and $\gcd(p,Q_j)=1$. Multiplication of both ciphertext components by $p^{-1}\bmod Q_j$ gives the phase
\[
T+p^{j-1}F\pmod{Q_j}.
\]
Interpret the result at plaintext modulus $p^{j-1}$, and verify its centered phase bound. This operation implements exact division only because the preceding plaintext is divisible by $p$; $p^{-1}$ is not an inverse in $\mathbb Z_{p^j}$. After $\varepsilon-r$ such divisions the plaintext modulus is $p^r$. There is no additional final division by $p^{\varepsilon-r}$.

\subsubsection{Discussion}

\para{The Reason for Modulus Switch from $q_l\rightarrow\hat q$:} The choice $\hat q\equiv1\pmod{p^\varepsilon}$ makes the overflow $\hat qK'$ reduce to $K'$ modulo the bootstrapping plaintext modulus $p^\varepsilon$. It does not by itself bound $K'$ or prove correctness. A valid BGV modulus switch must preserve the plaintext residue and leave the input error within the bound required above. The general switching operation need not require $\hat q$ to be a factor of the original modulus.

$ $

\para{The Choice of $\varepsilon$:} Select $\varepsilon$ from a deterministic input bound or an explicit probabilistic noise and secret-key model, together with a target failure probability. Increasing $\varepsilon$ also increases the term $p^\varepsilon\|E\|_\infty/q$, so its effect cannot be assessed by the digit gap alone. Larger $\varepsilon-r$ generally makes digit extraction more expensive. The cited paper gives the relevant parameter and failure analysis; no unconditional choice follows from the ring dimension alone.

$ $

\para{Generalization to $\Delta=p^r$:} The derivation above already includes every positive integer $r$, with $\varepsilon>r$. The scaling is $p^{\varepsilon-r}$, the rounding gap is $p^{\varepsilon-r}/2$, and digit extraction decreases the plaintext modulus from $p^\varepsilon$ to $p^r$. The BFV variant in \autoref{subsubsec:bfv-bootstrapping-summary} uses a different public preprocessing step; the common rounded-decryption structure does not make their ciphertext scaling rules interchangeable.







\subsubsection{Comparing the Bootstrapping in BFV, BGV, and CKKS}
\label{subsubsec:bootstrapping-differences}

In the BGV and BFV construction of~\cite{cryptoeprint:2022/1363}, public preprocessing is followed by a homomorphic inner product, linear coefficient/slot transformations, and exact digit extraction under a changing plaintext modulus. Correctness depends on an explicit decryption bound and the homomorphic circuit's error budget. BGV's intermediate exact divisions by $p$ reduce the plaintext modulus as well as the divisible phase; they are not generic inverses in the plaintext ring. BFV has a different scale-and-round preprocessing, as discussed in \autoref{subsubsec:bfv-bootstrapping-summary}.

CKKS bootstrapping lifts ciphertext coefficients, moves coefficient information into slots, approximately evaluates a periodic modular-reduction function, and applies the inverse linear transform. Its output retains the input approximation error and adds the circuit's approximation and encryption errors. A larger ciphertext modulus restores computational headroom, but does not in general recover precision already lost. Improved precision procedures require their own analysis; neither an unavoidable monotone error increase on every operation nor a universal precision reset follows from these outlines.
